\chapter{What a report must say about what was not screened}
\label{ch:reporting}

% Author's note (verification and reproduction): the four guidelines were
% retrieved as known items (OpenAlex lookups, 2026-11-23) and read in full
% from their publishers' copies; the commentary was returned by the first
% question's queries.  Counter-searches: q1y and q2z (see their COUNTER
% fields).  Sources: doc/bibliography.bib, keys Page2021 Rethlefsen2021
% Garritty2024 Stevens2024 Callaghan2024.

The last step of \cref{ch:procedure} writes down what was searched, what
was screened and what was not.
This chapter asks what the reporting guidelines require of a review that
does not screen everything: \emph{what must a report say about limits on
the search, about a partial screen, and about who or what screened?}
Unlike the four chapters before it, the question is about what documents
say, so the guidelines were retrieved as known items and read in full:
the reporting statement for systematic reviews, its extension for
literature searches, and the methods and reporting guidance for rapid
reviews.
The counter queries of the first and second questions
(\cref{ch:lit-search}) returned no work that contradicts what the
guidelines say, and one commentary on what they leave out, cited below.

\section{What the guidelines ask}
\label{sec:reporting-guidelines}

\begingroup\sloppy
\begin{description}
\item[Limits on the search] Item 7 of the statement for systematic
  reviews reads: \textcquote[Table~1, item~7]{Page2021}{Present the
  full search strategies for all databases, registers and websites,
  including any filters and limits used.}
  Its extension for searches asks for the reason too, in its item 9:
  \textcquote[Table~1, item~9]{Rethlefsen2021}{Specify that no limits were
  used, or describe any limits or restrictions applied to a search (e.g.,
  date or time period, language, study design) and provide justification
  for their use.}
  It also asks for the date of each search (item 13), the number of
  records each database returned (item 15) and how duplicates were
  removed (item 16).
\item[A partial screen] Of web search engines and websites, which can
  return more results than a team can screen, the extension says:
  \textcquote[p.~180]{Rethlefsen2021}{If review teams choose to review a
  limited set of results, it should be noted in the text, along with the
  rationale.}
  The reporting guidance for rapid reviews asks authors to
  \textcquote[Table, \enquote{Distinguishing the RR from an SR}]{Stevens2024}{\textins*{e}xplicitly
  describe why the product is an RR, noting the specific steps of conduct
  or methods that characterise the distinction from an SR}.
\item[Who or what screened] Item 8 of the statement reads:
  \textcquote[Table~1, item~8]{Page2021}{Specify the methods
  used to decide whether a study met the inclusion criteria of the review,
  including how many reviewers screened each record and each report
  retrieved, whether they worked independently, and if applicable, details
  of automation tools used in the process.}
  Its flow diagram asks, where automation tools were used, how many
  records a human excluded and how many the tools excluded.
\item[Two screeners, or one] The Cochrane guidance for rapid reviews
  recommends searching at least two bibliographic databases
  (recommendation 5) and, in recommendation 9, that two reviewers screen
  at least 20\,\% of the records, check their agreement and discuss the
  discrepancies:
  \textcquote[p.~7]{Garritty2024}{Teams can proceed with single screening
  if agreement is high (eg, $\kappa$ is $\geq$0.8).}
  It adds that the 20\,\% is \textcquote[p.~7]{Garritty2024}{based on our
  experience and is therefore not evidence based}.
\end{description}
\endgroup

\section{What they leave out}
\label{sec:reporting-gaps}

No guideline found asks for the reason a screen was stopped.
A commentary on stopping criteria notes of the statement for systematic
reviews that \textcquote[p.~2]{Callaghan2024}{it only stipulates reporting
the \enquote{details of automation tools used in the process}, without
requiring a rationale for stopping decisions}.
The procedure's account of the tier, the unscreened rest and the sample
therefore goes beyond what the guidelines require; it follows the
extension's request that a limited set of results be stated with its
rationale, and applies it to every ranked answer, not only to web
search engines.

\section{What the procedure writes}
\label{sec:reporting-procedure}

Taken together, the method of a search made by \cref{ch:procedure} says:
the databases searched and those that failed, each query verbatim with
its date and its hits per database (\texttt{scholar sessions show}), that
no limit was used, how duplicates were merged; the first tier's
definition and size, the records left unscreened, and the sample of them
with its seed, size, bearing count, rate and interval (\texttt{scholar
sessions tier-sample -{}-report}); who screened, by kind (a person, a
classifier with its model, an agent) and how many of the machine's
verdicts a person reviewed; and, where a second screen was made, the
agreement (\texttt{scholar sessions agreement}).
A negative result is worded as what was screened shows: none among the
records screened, with the sample's interval for the rest.

\section{Conclusion}
\label{sec:reporting-conclusion}

The guidelines ask that limits on a search be stated and justified, that
a limited set of ranked results be noted with its rationale, that the
number of records per database and the dates be given, and that the
report say how many reviewers screened each record and what automation
did; the rapid-review guidance asks for a second screener on a part of
the records and an agreement check, and states that its proportion is
not evidence based.
None asks for the reason a screen stopped, which the procedure states
anyway.
